\section{Introduction}
\label{sec:introduction}

\subsection{Background and Motivation}
Image classification maps an image to a semantic category. In an online product
catalogue, for example, an automatic classifier could help organize incoming
article images into clothing and accessory categories. The classification itself
is only one part of the system: the reliability of the result also depends on
the dataset, the representation of each image, the separation between training
and evaluation data, and the way model quality and computational cost are measured.

Fashion-MNIST is a useful setting for studying this pipeline. Its compact images
allow several neural network configurations to be investigated on the available
desktop GPU. At the same time, the recognition problem contains visually similar
categories and variation within each category. This combination supports a
practical assignment in which the data can be inspected carefully and the
experiments can be reproduced. The dataset was introduced by Xiao, Rasul, and
Vollgraf as a benchmark based on fashion article images\cite{xiao2017}.

Different neural architectures organize the same input information differently.
A linear classifier combines pixel values directly, while an MLP adds learned
nonlinear transformations. A CNN processes neighborhoods on the image grid.
An LSTM can read the image as an ordered sequence of rows, and a Transformer
can operate on a sequence of small image patches. Comparing these configurations
provides a way to investigate how spatial structure, sequence ordering, and
model capacity interact with a fixed classification task.

The motivation is therefore both predictive and experimental. A useful study
should identify which implemented configuration performs well, describe the
classes it confuses, and explain the resources it requires. A common data
protocol is essential: otherwise differences in splitting, normalization,
checkpoint selection, or hardware could obscure the comparison.

\subsection{Problem Statement}
The task is supervised, single-label, ten-class image classification. Let
\(X\in\{0,\ldots,255\}^{28\times28}\) be a raw grayscale image and
\(y\in\{0,\ldots,9\}\) its reference label. A classifier with parameters
\(\theta\) receives a preprocessed representation \(T(X)\) and produces
ten real-valued class scores, or logits:
\begin{equation}
 f_\theta(T(X))=\boldsymbol{z}\in\mathbb{R}^{10},
 \qquad \widehat{y}=\underset{k\in\{0,\ldots,9\}}{\arg\max}\;z_k.
 \label{eq:task}
\end{equation}
The output is one category for the complete image. Object detection,
segmentation, and multilabel classification are outside the scope of this task.

The main visual challenge is that 28 by 28 grayscale images retain only limited
detail. A shirt, pullover, coat, and T-shirt/top can have similar overall shapes,
while sleeve lengths, openings, and outlines vary within the same label. The
EDA in Section~\ref{sec:dataset} examines these properties directly. Such
observations motivate later hypotheses about confusion, rather than proving
which classes will be hardest for a trained classifier.

The experimental problem is to make the comparison interpretable. A difference
in accuracy can reflect input layout, capacity, regularization, optimization,
or the training budget. A difference in recorded training time can also reflect
the execution platform. The report must therefore distinguish measured behavior
of these particular configurations from general claims about architecture families.

\subsection{Objectives and Scope}
The study has four objectives:
\begin{enumerate}
 \item Establish a reproducible Fashion-MNIST data protocol, with a fixed
 stratified split, training-only normalization, and auditable data provenance.
 \item Compare Linear and MLP baselines with the implemented CNN, LSTM, and
 Transformer, keeping the dataset and nominal optimization budget consistent.
 \item Evaluate predictive quality using accuracy, macro-F1, class-specific
 measures, and confusion matrices, together with learning curves and image-level
 error examples.
 \item Measure parameter count, training cost, inference latency, and GPU memory
 on a documented reference platform, so that quality and efficiency can be
 interpreted together.
\end{enumerate}

This first stage completes the introduction and dataset analysis. It verifies
the local files, the existing split, class counts, pixel statistics, image
variation, exact duplicates, and preprocessing layouts. It does not introduce
new model scores or an architecture ranking. The remaining sections will
document the implementation, verify experimental artifacts, and interpret the
results under the agreed protocol.

The dataset is limited to Fashion-MNIST. No additional product photographs,
external training data, pretrained models, or synthetic images are introduced.
The shared configuration uses seed 42, ten epochs, batch size 128, Adam with a
learning rate of 0.001, and no data augmentation. These are the existing
assignment settings, not the outcome of a new hyperparameter search.

The local computational reference is the NVIDIA GeForce GTX 1660 SUPER with
6 GiB of GPU memory, using the existing Windows Python and CUDA PyTorch
environment. Existing Linear and MLP training summaries identify MPS as their
original training device; those durations cannot serve as a matched CUDA
training comparison. A subsequent experimental stage must obtain measurements
on this PC under consistent timing scopes and preserve the historical artifacts.
Equal epochs also do not imply equal wall-clock computation or equal parameter
capacity. The final discussion must retain these qualifications.

\subsection{Proposed Approach}
The approach begins by separating data inspection from model selection.
First, verify the local Fashion-MNIST files and reproduce the cached seed-42
training/validation indices. Next, inspect training images and compute class
balance, brightness distributions, nonzero-pixel coverage, spatial class
statistics, and exact-duplicate checks. Validation and test images do not
contribute to the normalization statistics or training-set visual summaries.

After inspection, scale each raw image to the unit interval and standardize it
using one mean and standard deviation estimated from the 54,000 training images.
Change only the layout for the selected model: a 784-value vector for Linear
and MLP, a one-channel spatial grid for CNN, 28 ordered row steps for LSTM, or
49 ordered 4 by 4 patches for Transformer. Section~\ref{sec:preprocessing}
defines these transformations and checks that all 784 pixel values are retained.

The later training stage will use the shared optimization protocol and choose
each model's checkpoint by its lowest validation cross-entropy. The fixed
checkpoint will then be evaluated on the official test partition. Results
will be interpreted through aggregate metrics, class-specific errors, learning
curves, and matched CUDA resource measurements. Additional runs, if performed,
will be identified explicitly as repetitions or ablations rather than mixed
with the original experiment.

This workflow makes each report claim traceable to code or an exported
artifact. In the present stage, the numerical EDA values, figure files, CSV
tables, and verification record are generated by
\artifact{report/generate_eda.py}. The next-stage guide defines the remaining
sections and the evidence needed before the report can be described as complete.
